% =====================================================================
\section{Programme home and the open selection layer}\label{sec:programme_home}
% =====================================================================

This section positions $\Cph$ within the broader cascade programme
and identifies what the operator does not deliver. The framing
matters for the hostile-review reading: $\Cph$ is the
\emph{response} primitive; \emph{selection} (which response
configuration the system picks under dynamics) is a separate layer
that this paper does not close.

\subsection{Programme position}

$\Cph$ is one instance of a polynomial-in-$L$ functional on a
finite-dimensional substrate. The published companion adaptive-
closure-transport
preprint~\citep{SmartAdaptiveClosureTransport2026} positions a
homeostatic regulariser $R_{\mathrm{hom}}$ in the 4-tuple
$(M, L_M, W, R_{\mathrm{hom}})$ as a member of the same
polynomial-in-$L$ family. We use this as the only published
programme-level reference point for $\Cph$. We make no
classification claim, no family-membership theorem, and no
forecast about any other members of this programme; any
broader programme arc lies outside the operator-witness scope of
this paper.

\subsection{Response vs selection}

The closure response $\psi = \Cph^{-1} f$ is determined by the
chosen substrate plus the design-level shift: $\Cph$ is fixed by
the substrate $\Rsixhundred$ and the design-level choice
$\Ph^{-2}$, and the response is the resulting linear inverse.
This is a \emph{response} primitive. It does \emph{not} answer:
\begin{itemize}\itemsep=2pt
\item Why this substrate? (Selection across regular 4-polytopes
  $\{24\text{-cell}, 600\text{-cell}, 120\text{-cell}\}$.)
\item Why this shift? (Selection of $\Ph^{-2}$ over an arbitrary
  positive constant.)
\item How does the system pick a response configuration over
  time? (Crystallisation / Lyapunov descent dynamics on a
  $W$-trajectory.)
\end{itemize}

The selection layer is open. The companion adaptive-closure-
transport preprint~\citep{SmartAdaptiveClosureTransport2026}
proposes a Lyapunov $V(W)$ on the reduced flow, an edge-space
decomposition under the $2I$ binary-icosahedral symmetry of
$\Rsixhundred$, and a full reduced-flow convergence theorem on
$W$-trajectories; that paper explicitly does \emph{not} deliver
those, and we do not deliver them here. The aria-chess
companion~\citep{SmartAriaChess2026} likewise stays inside
substrate-witness scope and does \emph{not} deliver a selection
theorem. The two empirical witnesses landed in this paper provide
external consistency checks on the \emph{response} primitive
without reducing or addressing the selection gap.

\subsection{What this paper closes vs leaves open}

\paragraph{Closes (at the operator-witness level).}
\begin{itemize}\itemsep=2pt
\item The operator $\Cph$ is well-defined and positive definite
  on any $(M, L_M)$ satisfying (H1)--(H3); the operator-norm
  identity $\|\Cph^{-1}\| = \Ph^{2}$ holds whenever
  $\inf \sigma(L_M) = 0$ (e.g.\ on a connected finite graph with
  the standard combinatorial Laplacian, or on the full-line
  continuum case). On substrates where $\inf \sigma(L_M) > 0$
  (e.g.\ Dirichlet-boundary continuum cases) the bound
  $\|\Cph^{-1}\| \leq \Ph^{2}$ holds and is generally strict
  (\S\ref{sec:definition}).
\item The 600-cell instance $\Rsixhundred$ has the construction
  described (\S\ref{sec:substrate}) and the Laplacian spectrum of
  Table~\ref{tab:spectrum}, both reproduced numerically
  (\texttt{repro/verify\_kernel.py}).
\item Discrete-to-continuum agreement at per-vertex Pearson
  correlation $0.976$ on the unweighted variant, with the unweighted
  variant winning the geometry-only criterion against two
  $\Ph$-cocycle weighted controls (\S\ref{sec:agreement}).
\item Same fixed $\Cph$ on same fixed $\Rsixhundred$ appears as
  the shared response primitive underneath two domain-disjoint empirical works in
  qualitatively distinct regimes (\S\ref{sec:passive_witness},
  \S\ref{sec:active_witness}).
\end{itemize}

\paragraph{Leaves open.}
\begin{itemize}\itemsep=2pt
\item \emph{First-principles derivation of $\Ph^{-2}$.} Reported
  as a design-level shift; not derived from a closure functional
  or symmetry argument.
\item \emph{Substrate-uniqueness ablation.} The 600-cell choice is
  post-hoc motivated by the empirical landings; alternative regular
  4-polytopes are an explicit ablation build, not a discharged
  comparison.
\item \emph{Kernel-uniqueness on either empirical landing.} The
  b-anomaly free-width Gaussian alternative (fits comparably with
  one extra shape parameter), the AIC tie
  ($w_{\mathrm{VFD}}=0.348$ vs $w_{\mathrm{FREE\_C9}}=0.652$), and
  the Mode-B linearised-to-non-linear refit drift caveat are
  inherited verbatim from~\citep{SmartBAnomaly2026}.
\item \emph{Selection theorem on ACT.} Lyapunov $V(W)$, edge-space
  decomposition under $2I$-equivariance, full reduced-flow
  convergence — all explicitly not delivered
  in~\citep{SmartAdaptiveClosureTransport2026} and not delivered
  here.
\item \emph{Family-membership theorem.} The programme-home
  positioning of cascade Lyapunov functionals as members of the
  same polynomial-in-$L$ family is reported as
  \emph{programme-positioned}, not formally classified.
\end{itemize}
